% from ega4-2-fr.tex
\oldpage[IV-2]{72}

l'on suppose en outre \emph{géométriquement irréductible}~; si de plus $X$ est
localement noethérien, il y a alors un voisinage ouvert de $x$ dans $X$ qui est
\emph{géométriquement intègre}.
\end{env}

Rappelons que si un préschéma localement noethérien est ponctuellement intègre,
il est localement intègre (I, 6.1.13). Si un $k$-préschéma $X$, localement de
type fini sur $k$, est géométriquement ponctuellement intègre sur $k$, il
résulte donc de cette remarque que, pour toute extension $k'$ de $k$,
$X\otimes_k k'$ est \emph{localement intègre} (on dit encore dans ce cas que
$X$ est \emph{géométriquement localement intègre}).

\begin{proposition}[4.6.15]
\label{IV.2.4.6.15-fr}
\begin{enumerate}
\item[(i)] Soient $k$ un corps~; $X$ un $k$-préschéma de type fini. Pour que $X$
soit géométriquement ponctuellement intègre, il faut et il suffit qu'il soit
géométriquement réduit et que le nombre géométrique de ses composantes
irréductibles soit égal au nombre géométrique de ses composantes connexes.
\item[(ii)] Soit $X$ un préschéma localement de type fini sur $k$. Si $X$ est
géométriquement ponctuellement intègre en un point $x$, il existe un voisinage
ouvert $U$ de $x$ qui est géométriquement ponctuellement intègre. En d'autres
termes, l'ensemble des $x\in X$ où $X$ est géométriquement ponctuellement
intègre est ouvert dans $X$.
\end{enumerate}
\end{proposition}

\begin{proof}
\begin{enumerate}
\item[(i)] Étant donnée une extension algébriquement close $\Omega$ de $k$,
$X_{(\Omega)}$ est ponctuellement intègre (ou localement intègre, ce qui
revient au même puisque $X_{(\Omega)}$ est noethérien) si et seulement s'il est
réduit et si le nombre de ses composantes connexes est égal au nombre de ses
composantes irréductibles (I, 6.1.10). Il suffit alors d'appliquer (4.5.1),
(4.6.1) et la première remarque de (4.6.14).

\item[(ii)] La question étant locale sur $X$, on peut se borner au cas où $X$
est un préschéma de type fini sur $k$. En outre, on sait qu'il y a un voisinage
ouvert de $x$ dans $X$ qui est géométriquement réduit (4.6.13), donc on peut
supposer $X$ géométriquement réduit. Il existe alors une extension finie $k'$
de $k$ tel que $X'=X\otimes_k k'$ soit réduit et que ses composantes
irréductibles soient géométriquement irréductibles (4.6.8). Soit
$p:X'\to X$ la projection canonique, qui est un morphisme fini surjectif, et
soient $x'_j$ ($1\leqslant j\leqslant r$) les points de $p^{-1}(x)$~; par
hypothèse, $X'$ est intègre en chacun des $x'_j$, donc chaque $x'_j$ a un
voisinage ouvert $V'_j$ dans $X'$ qui est intègre. Comme $p$ est fermé
(II, 6.1.10), chacun des $p(V'_j)$ est un voisinage de $x$, donc il existe un
voisinage ouvert $U$ de $x$ tel que $p^{-1}(U)$ soit contenu dans la réunion
des $V'_j$, et soit par suite intègre en chacun de ses points, donc localement
intègre. Les composantes irréductibles de $p^{-1}(U)$ sont donc ouvertes et
deux à deux disjointes, et comme elles sont géométriquement irréductibles
(4.5.9), $p^{-1}(U)=U\otimes_k k'$ est géométriquement ponctuellement intègre
en vertu de (i), et il en est donc de même de $U$ par définition.
\end{enumerate}
\end{proof}

\begin{proposition}[4.6.16]
\label{IV.2.4.6.16-fr}
Soient $k$ un corps, $X$ un $k$-préschéma localement noethérien, $\mathscr F$
un $\mathcal O_X$-Module cohérent. Les conditions suivantes sont
équivalentes~:
\begin{enumerate}
\item[a)] Pour toute extension de type fini $k'$ de $k$, si l'on pose
$X'=X\otimes_k k'$, alors $\mathscr F'=\mathscr F\otimes_k k'$ est un
$\mathcal O_{X'}$-Module réduit (3.2.2).
\item[b)] Pour toute extension finie radicielle $k'$ de $k$, $\mathscr F'$ est
un $\mathcal O_{X'}$-Module réduit.
\item[c)] $\mathscr F$ est réduit, et si $\mathscr J$ est l'Idéal de
$\mathcal O_X$ annulateur de $\mathscr F$, le sous-préschéma fermé défini par
$\mathscr J$ est géométriquement réduit sur $k$.
\end{enumerate}

